\heading{数学物理方法（上）-第8次作业}

\begin{question}
将下列函数在指定点展开为Taylor级数，并给出其收敛半径：
\begin{enumerate}
  \item \(\arctan z\) 的主值，在 \(z=0\) 处展开；

  \item \(\ln \frac{1+z}{1-z}\) 在 \(z=\infty\) 处展开，规定 \(\bigl. \ln \frac{1+z}{1-z} \bigr|_{z=\infty}=(2k+1) \pi \ii\)。
\end{enumerate}
\begin{solution}
\begin{enumerate}
  \item 有 \(\arctan z = \frac{1}{2\ii} \ln \frac{1+\ii z}{1-\ii z}\)，在 \(z=0\) 处展开，则有
\[
\arctan z = \frac{1}{2\ii} \ln\!\left(1+\frac{2\ii z}{1-\ii z}\right)
= \frac{1}{2\ii} \sum_{n=1}^\infty \frac{-1}{n} \left(\frac{-2\ii z}{1-\ii z}\right)^n
= \frac{1}{2\ii} \sum_{n=1}^\infty \frac{-1}{n} (-2\ii z)^n \sum_{m=0}^\infty \binom{-n}{m} (-\ii z)^m
\]
\[
= \frac{-1}{2\ii} \sum_{n=1}^\infty \sum_{m=0}^\infty \frac{2^n}{n} \binom{-n}{m} (-\ii z)^{m+n}
\]
令 \(k = m+n\)，则有 \(n = k-m\)，代入有
\[
\arctan z = \frac{-1}{2\ii} \sum_{k=1}^\infty \sum_{m=0}^{k-1} \frac{2^{k-m}}{k-m} \binom{m-k}{m} (-\ii z)^k
= \sum_{k=1}^\infty \frac{-1}{2\ii} \left[\sum_{m=0}^{k-1} \frac{2^{k-m}}{k-m} \binom{m-k}{m} \ii^{-k}\right] z^k
\]

设 \(S_k = \sum_{m=0}^{k-1} \frac{2^{k-m}}{k-m} \binom{m-k}{m}\)，有
\[
\binom{m-k}{m}=(-1)^{k-m} \binom{k-m +m-1}{m}=(-1)^{k-m} \binom{k-1}{m}
\]

因此可以得到，
\[
S_k = \sum_{m=0}^{k-1} \frac{(-2)^{k-m}}{k-m} \binom{k-1}{m}
= \sum_{m=1}^k \frac{(-2)^m}{m} \binom{k-1}{k-m}
= \sum_{m=1}^k \frac{(-2)^m}{m} \binom{k-1}{m-1}
\]

注意到，
\[
\dv{}{x} \sum_{j=1}^n \frac{x^j}{j} \binom{n-1}{j-1}
= \sum_{j=1}^n x^{j-1} \binom{n-1}{j-1} = (1+x)^{n-1}
\]

因此有，
\[
\sum_{j=1}^n \frac{x^j}{j} \binom{n-1}{j-1} = \frac{(1+x)^n-1}{n}
\]

所以可以得到，
\[
S_k = \frac{(-1)^k-1}{k} = \begin{cases}
0, & k \text{ 为 偶 数},\\
-\frac{2}{k}, & k \text{ 为 奇 数}
\end{cases}
\]

\[
\arctan z = \sum_{k=1}^\infty \frac{-\ii^{-k}}{2\ii} S_k z^k
= \sum_{k=1}^\infty \frac{(-1)^{k-1} z^{2k-1}}{2k-1}
\]

其收敛半径为 1。(因为 \(\arctan z\) 在 \(z=\pm \ii\) 处为奇点。)

\item 代入 \(z=1/t\)，有 \(\ln \frac{1+z}{1-z} = \ln \frac{t+1}{t-1}\)，由题 \(\ln (-1+z)=(2k+1) \pi \ii + \ln(1-z)=(2k+1)\pi \ii - \sum_{k=1}^\infty \frac{z^k}{k}\)。因此可以得到
\[
\ln \frac{t+1}{t-1} = (2k+1)\pi \ii + \ln\!\left(\frac{1+t}{1-t}\right)
= (2k+1)\pi \ii + \operatorname{v.p.}\{\ln(1+t) - \ln(1-t)\}
\]
\[
= (2k+1) \pi \ii + \sum_{k=1}^\infty \frac{t^k-(-t)^k}{k}
= (2k+1)\pi \ii + \sum_{k=1}^\infty \frac{2t^{2k-1}}{2k-1}
\]

因此可以得到
\[
\ln \frac{1+z}{1-z} = (2k+1)\pi \ii + \sum_{k=1}^\infty \frac{2}{2k-1} \left(\frac{1}{z}\right)^{2k-1}
\]

函数 \(\ln \frac{1+z}{1-z}\) 的奇点在 \(z=\pm 1\)，因此上述表达式收敛域为 \(\{z:|z|>1\}\)。
\end{enumerate}
\end{solution}
\end{question}

\begin{question}
求下列无穷级数之和，注意给出相应的收敛区域：
\[
\sum_{n=0}^\infty \sum_{m=0}^\infty \sum_{k=0}^\infty \frac{(n+m+k)!}{n!\,m!\,k!} \left(\frac{z}{3}\right)^{m+n+k}
\]
\begin{solution}
利用多项式展开，有
\[
(a+b+c)^s = \sum_{m+n+k=s} \frac{s!}{m!\,n!\,k!} a^m b^n c^k
\]

因此上述求和可以表示为
\[
S = \sum_{s=0}^\infty \left(\frac{z}{3}+\frac{z}{3}+\frac{z}{3}\right)^s
= \sum_{s=0}^\infty z^s = \frac{1}{1-z}
\]

收敛区域为 \(\{z:|z|<1\}\)。
\end{solution}
\end{question}

\begin{question}
求下列函数（取主值分支）在 \(z=0\) 附近的级数展开：
\begin{enumerate}
  \item \(-\ln(1-z) \ln(1+z)\);

  \item \(\ln(1+z^2) \arctan z\)。
\end{enumerate}
\begin{solution}
\begin{enumerate}
  \item
\[
-\ln(1-z) \ln(1+z) = \sum_{k=1}^\infty \frac{z^k}{k} \sum_{l=1}^\infty -\frac{(-z)^l}{l}
\]

令 \(n = k+l\)，则有
\[
-\ln(1-z) \ln(1+z) = \sum_{n=2}^\infty \left[\sum_{k=1}^{n-1} \frac{(-1)^{n-k+1}}{k(n-k)}\right] z^n
\]

由对称性分析可知，当 \(n\) 为奇数时，\(k=\frac{n}{2} \pm (\frac12 + q)\) 时，\((-1)^{n-k_++1}=(-1)^{n/2+1/2-q}\)，\((-1)^{n-k_-+1}=(-1)^{n/2-1/2+q}\)，两者次数相差 \(2q-1\)，因此
\[
(-1)^{n-k_++1} \frac{1}{k_+(n-k_+)} + (-1)^{n-k_-+1} \frac{1}{k_-(n-k_-)} = 0
\]

因此有
\[
-\ln(1-z) \ln(1+z) = \sum_{n=1}^\infty \left[\sum_{k=1}^{2n-1} \frac{(-1)^{k+1}}{k(2n-k)}\right] z^{2n}
\]

  \item
\[
\ln(1+z^2) \arctan z
= \sum_{k=1}^\infty -\frac{(-z^2)^k}{k} \sum_{l=1}^\infty \frac{(-1)^{l-1} z^{2l-1}}{2l-1}
\]

令 \(n = k+l\)，可以得到
\[
\ln(1+z^2) \arctan z
= \sum_{n=2}^\infty \left[\sum_{k=1}^{n-1} \frac{(-1)^n}{k\,(2(n-k)-1)}\right] z^{2n-1}
\]
\end{enumerate}
\end{solution}
\end{question}
